\section{Formal Mathematical Proofs}

\begin{theorem}[Bitwise Invertibility of Residue-Hiding Encryption]
Let $T \in \mathbb{R}^{B \times C \times D}$ be a floating-point tensor represented in IEEE 754 double precision format. The Residue-Hiding encryption transformation $\mathcal{E}_i(T) = (T + \delta_i) \phi^{-(i \bmod 13)}$ and decryption operator $\mathcal{D}_i(\hat{T}) = \hat{T} \phi^{i \bmod 13} - \delta_i$ form a bi-Lipschitz homeomorphism with numerical reconstruction error bounded by $\epsilon_{\text{mach}} \approx 2.22 \times 10^{-16}$.
\end{theorem}

\begin{proof}
Consider the decryption error $E = |\mathcal{D}_i(\mathcal{E}_i(T)) - T|$.
Under floating-point arithmetic with machine epsilon $\epsilon$, addition and scalar multiplication satisfy:
\begin{equation}
    \text{fl}(T + \delta_i) = (T + \delta_i)(1 + \epsilon_1), \quad |\epsilon_1| \le \epsilon
\end{equation}
\begin{equation}
    \text{fl}\left(\text{fl}(T + \delta_i) \cdot \phi^{-n}\right) = (T + \delta_i)\phi^{-n}(1 + \epsilon_1)(1 + \epsilon_2)
\end{equation}
Multiplying by $\phi^n$ in exact arithmetic yields:
\begin{equation}
    (T + \delta_i)(1 + \epsilon_1)(1 + \epsilon_2) - \delta_i = T + (T + \delta_i)(\epsilon_1 + \epsilon_2 + \epsilon_1\epsilon_2)
\end{equation}
Subtracting $T$ and taking the norm yields $\|E\| \le (\|T\| + \delta_i) \cdot 2\epsilon + \mathcal{O}(\epsilon^2)$. Since $\|T\| \le M$ and $\delta_i \in [0, 1)$, $\|E\| = \mathcal{O}(\epsilon_{\text{mach}})$. Hence, QSV encryption guarantees loss-free retrieval up to double-precision machine epsilon.
\end{proof}

\begin{theorem}[KV-Cache Memory Footprint Reduction]
Let $L$ be the total sequence length processed, and let $C$ be the shard capacity threshold ($C = 1024$). Under $k$-step $\phi^\infty$ Shard Folding compression ratio $\rho = \phi^{-6k}$, the total memory footprint $M(L)$ satisfies:
\begin{equation}
    M(L) \le 2 \cdot B \cdot D \cdot \left((L \bmod C) + \phi^{-6k} \cdot \left\lfloor \frac{L}{C} \right\rfloor \cdot C\right)
\end{equation}
\end{theorem}

\begin{proof}
The uncompressed active tokens consume at most $(L \bmod C)$ sequence slots. The number of compressed historical shards is $S = \lfloor L / C \rfloor$. Each shard of length $C$ is compressed by factor $\rho = \phi^{-6k}$. Summing memory allocations:
\begin{equation}
    M(L) = 2BD \left(L_{\text{active}} + \sum_{j=1}^S C_{\text{compressed}}\right) = 2BD \left((L \bmod C) + S \cdot C \cdot \phi^{-6k}\right)
\end{equation}
For $k = 3$, $\phi^{-18} = \frac{1}{(1.61803398875)^{18}} \approx 0.000173$. Memory overhead for historical sequence blocks is reduced by over $99.98\%$, establishing near-constant memory complexity.
\end{proof}
